\documentclass[12pt]{article}
\usepackage[utf8]{inputenc}
\usepackage[T1]{fontenc}
\usepackage{hyperref}
\title{Xukur\'u Language Resources for K-12 Curriculum Development}
\author{Compiled by Claude Sonnet 5.5 (claude-sonnet-5-5) for Andr\'as Kornai}
\date{}
\begin{document}
\maketitle
\section{Introduction}
Xukur\'u (ISO 639-3 xoo; Glottocode xuku1239) is a dormant language of Brazil with no remaining first-language speakers. It is known only from a few word lists. Search results showed no school curriculum in the language. The Xukuru do Orub\'a people of Pernambuco are working on revitalization, so any curriculum would have to build on those community efforts. Little exists, and this report does not pad it.

\section{Primary Educational Resources}
A 2025 article in Revista da ABRALIN describes an illustrated dictionary of about 94 words of the language, made as a revitalization and identity tool. The dictionary itself was not seen in the search results, only the article about it. \url{https://revista.abralin.org/index.php/abralin/article/download/2363/2971/19332}

\section{Linguistic and Reference Materials}
\begin{itemize}
\item Wikipedia, Xukuru language (overview; the language is poorly documented, extinct, and recorded in a few word lists): \url{https://en.wikipedia.org/wiki/Xukuru_language}
\item UFPE repository record for a dissertation by Daniel Bernardo Rocha Guimar\~aes de Souza, found in a search on the Xukuru. Its content was not checked. \url{https://repositorio.ufpe.br/handle/123456789/52167?mode=full}
\end{itemize}

\section{Religious and Community Resources}
An article on the Xukuru do Orub\'a people and their mobilizations for rights, based on memory, in the Agreste of Pernambuco: \url{https://www.e-publicacoes.uerj.br/revistaceaju/article/view/65134}

\section{Digital and Audio Resources}
No audio, video or software resources were found in the search results.

\section{Curriculum Development Context}
There are no native speakers and no known school instruction in the language. Any curriculum would be a revitalization programme built on the small documented vocabulary, together with the community's own initiatives (the Xukuru do Orub\'a). Developers should consult the community directly and should not assume that more material exists than was found here.
\end{document}